\documentclass{jecde}
\title{Designing reproducible engineering research: a guide to the JECDE article format}
\jecshorttitle{Article layout and author guidance}
\jecarticletype{Template demonstration}
\jecauthors{JECDE Editorial Design}
\jecaffiliations{Author template development edition; no editorial appointment is implied.}
\jeccorrespondence{Template v1.0\quad |\quad Compile with pdfLaTeX\quad |\quad 30 October 2026}
\jecabstract{This document demonstrates the proposed article design for the Journal of Electrical, Computer and Digital Engineering. It illustrates title hierarchy, a full-width abstract, two-column body text, numbered equations, vector figures, formal tables and an algorithm with a coloured header. The content is instructional rather than a published research article. Mathematical examples and graphics explain formatting conventions; they do not represent experimental findings. Publication identifiers and acceptance dates are deliberately absent because none have been assigned.}
\jeckeywords{electrical engineering; electronics; computing; communications; artificial intelligence; cybersecurity; digital technologies}
\begin{document}
\maketitle
\section{Purpose and scope}
This template provides a consistent visual language across the journal's seven subject tracks. Electrical engineering, electronics, computing, communications, artificial intelligence, cybersecurity and digital technologies share the same article structure while retaining discipline-specific methods and notation.



\begin{jecnote}{Design demonstration}
This is an author template, not a published article. No ISSN, DOI, volume, acceptance date or indexing status is claimed. Replace all instructional content before submission.
\end{jecnote}

\subsection{Structure and contribution}
An original research article should state a precise problem, explain its novelty in context, define a reproducible method and support its conclusions with evidence. A review article should describe its literature-selection method and provide critical synthesis.

Use numbered sections and concise subsection titles. Define acronyms at first appearance. Cite figures, tables and equations in the text before discussing their implications.

\section{Mathematics and notation}
Mathematical notation should be consistent throughout the manuscript. For example, a generic engineering allocation problem can be expressed as
\begin{align}
\underset{\boldsymbol{x}}{\operatorname{minimize}}\quad &J(\boldsymbol{x})=\sum_{i=1}^{N}w_i f_i(x_i),\label{eq:objective}\\
\text{subject to}\quad &\sum_{i=1}^{N}x_i\leq B,\quad x_i\geq0.\label{eq:constraints}
\end{align}
Here, $x_i$ denotes the resource allocated to component $i$, $B$ is the total budget, $w_i$ is a dimensionless importance weight and $f_i$ is a problem-specific cost function. Equations~\eqref{eq:objective} and~\eqref{eq:constraints} are formatting examples; a research manuscript must justify its own assumptions and units.

\subsection{Reproducibility}
Report software, data, parameters and evaluation conditions. Distinguish analytical, simulated and measured results. Table~\ref{tab:tracks} summarizes track-specific reporting checks.

\begin{table*}[t]
\caption{Recommended checks across the seven subject tracks.}\label{tab:tracks}
\centering\small
\begin{tabularx}{\textwidth}{@{}lXX@{}}\toprule
Track & Typical evidence & Reporting emphasis\\\midrule
Electrical engineering & Analysis, simulation or measurements & Units, operating conditions, stability and losses\\
Electronics & Circuit models and device measurements & Component values, fabrication and uncertainty\\
Computing & Benchmarks and software artifacts & Environment, complexity and reproducibility\\
Communications & Link or network evaluation & Channels, traffic, baselines and metrics\\
Artificial intelligence & Training and independent testing & Data splits, leakage prevention and robustness\\
Cybersecurity & Threat models and security evaluation & Adversary capabilities and ethical handling\\
Digital technologies & Integrated system evaluation & Architecture, deployment and interoperability\\\bottomrule
\end{tabularx}
\end{table*}

\section{Figures and visual evidence}
Use vector graphics where practical and ensure labels remain readable at the final column width. Figure~\ref{fig:workflow} demonstrates a compact workflow with consistent shapes and line weights.
\begin{figure}[ht]
\centering
\begin{tikzpicture}[font=\sffamily\small,>=stealth]
\tikzstyle{stage}=[draw=jecgreen,fill=jecsage,rounded corners=2pt,minimum width=5.2cm,minimum height=.65cm,align=center]
\node[stage](a)at(0,0){Research question};
\node[stage](b)at(0,-1){Model and assumptions};
\node[stage](c)at(0,-2){Evaluation and uncertainty};
\node[stage](d)at(0,-3){Evidence and interpretation};
\draw[->,jecgreen,thick](a)--(b);\draw[->,jecgreen,thick](b)--(c);\draw[->,jecgreen,thick](c)--(d);
\end{tikzpicture}
\caption{A research workflow illustrated with editable vector artwork.}\label{fig:workflow}
\end{figure}



\subsection{Tables}
Table~\ref{tab:components} uses formal horizontal rules rather than shaded cells or a dense grid. Keep units in column headings and explain abbreviations in the surrounding text.
\begin{table}[ht]
\centering\caption{Article components and their purpose.}\label{tab:components}
\small\begin{tabularx}{\columnwidth}{@{}lX@{}}\toprule
Component & Purpose\\\midrule
Abstract & Problem, approach and supported findings\\
Methods & Assumptions and reproducible procedure\\
Results & Evidence, comparisons and uncertainty\\
Discussion & Interpretation and limitations\\\bottomrule
\end{tabularx}
\end{table}

\section{Algorithms}
Use an explicit input/output structure and meaningful operations. The numbered example below demonstrates the journal's header-bar style. It is a manuscript-checking procedure, not a novel research algorithm.
\begin{jecalgorithm}{Manuscript readiness check}\label{alg:check}
\textbf{Input:} manuscript $M$, checklist $C$.\\
\textbf{Output:} unresolved issues $U$.
\begin{enumerate}[label=\arabic*.,leftmargin=1.4em]
\item Initialize $U\leftarrow\varnothing$.
\item For each requirement $c\in C$, inspect the relevant part of $M$.
\item If $c$ is not satisfied, append an evidence-based note to $U$.
\item Verify figures, equations, references and declarations.
\item Return $U$ for author revision.
\end{enumerate}
\end{jecalgorithm}



\section{Reporting and interpretation}
A results section should explain how the evidence answers the research question. Report the evaluation conditions and appropriate baselines, then discuss why changes occur. Include uncertainty or repeatability information where applicable.



\subsection{References}
References should be accurate, relevant and cited in order of first appearance. A sample citation to a typesetting reference is provided here~\cite{lamport}. Replace it with sources appropriate to the submitted research.

\section{Wide material and final checks}
Two-column figures and tables are supported by the standard \texttt{figure*} and \texttt{table*} environments. Place wide content near its first discussion and allow the typesetting engine to find a suitable position.





\section{Conclusion}
The JECDE format combines a recognizable opening page with conventional scholarly typesetting. Authors should preserve readability, provide inspectable evidence and ensure that presentation supports the scientific contribution.

\section*{Declarations}
\textbf{Funding:} No external funding supported this template demonstration.

\textbf{Competing interests:} No research competing interests are asserted in this design sample.

\textbf{Data and code availability:} This document contains no experimental dataset. Its editable LaTeX source accompanies the template package.

\textbf{Author contributions:} The document is a design and formatting demonstration, not a research authorship record.

\begin{thebibliography}{1}
\bibitem{lamport} L. Lamport, \emph{LaTeX: A Document Preparation System}, 2nd ed. Addison-Wesley, 1994.
\end{thebibliography}
\end{document}
